% ===== VARIANT: QUANTITATIVE ASSET MANAGEMENT =======================
% Angle: portfolio construction, index replication, liquidity, ESG,
% asset allocation and performance. Pricing kept but subordinated.
% ===================================================================
% =====================================================================
%  Shared preamble for all CV variants.
%  Each variant does:   % =====================================================================
%  Shared preamble for all CV variants.
%  Each variant does:   % =====================================================================
%  Shared preamble for all CV variants.
%  Each variant does:   \input{cv-preamble.tex}
%  Do not put content here — only packages, formatting and macros.
% =====================================================================
\documentclass[letterpaper,11pt]{article}

\usepackage{latexsym}
\usepackage[empty]{fullpage}
\usepackage{titlesec}
\usepackage{marvosym}
\usepackage[usenames,dvipsnames]{color}
\usepackage{verbatim}
\usepackage{enumitem}
\usepackage[hidelinks]{hyperref}
\usepackage{fancyhdr}
\usepackage[english]{babel}
\usepackage{tabularx}
\usepackage{fontawesome5}
\usepackage{multicol}
\usepackage{graphicx}
\setlength{\multicolsep}{-3.0pt}
\setlength{\columnsep}{-1pt}
\input{glyphtounicode}
\usepackage{paracol}

\pagestyle{fancy}
\fancyhf{}
\fancyfoot{}
\renewcommand{\headrulewidth}{0pt}
\renewcommand{\footrulewidth}{0pt}

\addtolength{\oddsidemargin}{-0.5in}
\addtolength{\evensidemargin}{-0.5in}
\addtolength{\textwidth}{0.9in}
\addtolength{\topmargin}{-.3in}
\addtolength{\textheight}{0.9in}
\hypersetup{
    colorlinks=true,
    linkcolor=blue,
    urlcolor=blue,
    filecolor=blue,
    pdfpagemode=FullScreen
}
\urlstyle{same}

\raggedbottom
\raggedright
\setlength{\tabcolsep}{0in}

\titleformat{\section}{
  \vspace{-4pt}\scshape\raggedright\large\bfseries
}{}{0em}{}[\color{black}\titlerule \vspace{-5pt}]

\pdfgentounicode=1

\newcommand{\resumeItem}[1]{
  \item\small{{#1 \vspace{-2pt}}}
}

\newcommand{\resumeSubheading}[4]{
  \vspace{-2pt}\item
    \begin{tabular*}{1.0\textwidth}[t]{l@{\extracolsep{\fill}}r}
      \textbf{#1} & \textbf{\small #2} \\
      \textit{\small#3} & \textit{\small #4} \\
    \end{tabular*}\vspace{-7pt}
}

\newcommand{\resumeProjectHeading}[2]{
    \item
    \begin{tabular*}{0.97\textwidth}{l@{\extracolsep{\fill}}r}
      \small#1 & \textbf{\small #2}\\
    \end{tabular*}\vspace{-7pt}
}

\renewcommand\labelitemi{$\vcenter{\hbox{\tiny$\bullet$}}$}
\renewcommand\labelitemii{$\vcenter{\hbox{\tiny$\bullet$}}$}

\newcommand{\resumeSubHeadingListStart}{\begin{itemize}[leftmargin=0.0in, label={}]}
\newcommand{\resumeSubHeadingListEnd}{\end{itemize}}
\newcommand{\resumeItemListStart}{\begin{itemize}}
\newcommand{\resumeItemListEnd}{\end{itemize}\vspace{-5pt}}

% ---------------------------------------------------------------------
% \cvheader{positioning line}{focus line}
% ---------------------------------------------------------------------
\newcommand{\cvheader}[2]{%
\begin{center}
    {\Huge \scshape Franck Deba} \\[0.4em]
    {\small \textit{#1}} \\[0.25em]
    {\small #2}
\end{center}
\vspace{0.5pt}
\small
\raisebox{-0.1\height}{\faPhone}~+33 (0)6 05 77 47 51 \hspace{1em}
\raisebox{-0.2\height}{\faEnvelope}~\href{mailto:debafranck@gmail.com}{\underline{debafranck@gmail.com}} \hspace{1em}
\raisebox{-0.2\height}{\faLinkedin}~\href{https://linkedin.com/in/deba-franck}{\underline{deba-franck}} \hspace{1em}
\raisebox{-0.2\height}{\faGithub}~\href{https://github.com/usenameFD}{\underline{usenameFD}} \hspace{1em}
\raisebox{-0.2\height}{\faGlobe}~\href{https://usenamefd.github.io/}{\underline{usenamefd.github.io}}
\vspace{-8pt}
}

%  Do not put content here — only packages, formatting and macros.
% =====================================================================
\documentclass[letterpaper,11pt]{article}

\usepackage{latexsym}
\usepackage[empty]{fullpage}
\usepackage{titlesec}
\usepackage{marvosym}
\usepackage[usenames,dvipsnames]{color}
\usepackage{verbatim}
\usepackage{enumitem}
\usepackage[hidelinks]{hyperref}
\usepackage{fancyhdr}
\usepackage[english]{babel}
\usepackage{tabularx}
\usepackage{fontawesome5}
\usepackage{multicol}
\usepackage{graphicx}
\setlength{\multicolsep}{-3.0pt}
\setlength{\columnsep}{-1pt}
\input{glyphtounicode}
\usepackage{paracol}

\pagestyle{fancy}
\fancyhf{}
\fancyfoot{}
\renewcommand{\headrulewidth}{0pt}
\renewcommand{\footrulewidth}{0pt}

\addtolength{\oddsidemargin}{-0.5in}
\addtolength{\evensidemargin}{-0.5in}
\addtolength{\textwidth}{0.9in}
\addtolength{\topmargin}{-.3in}
\addtolength{\textheight}{0.9in}
\hypersetup{
    colorlinks=true,
    linkcolor=blue,
    urlcolor=blue,
    filecolor=blue,
    pdfpagemode=FullScreen
}
\urlstyle{same}

\raggedbottom
\raggedright
\setlength{\tabcolsep}{0in}

\titleformat{\section}{
  \vspace{-4pt}\scshape\raggedright\large\bfseries
}{}{0em}{}[\color{black}\titlerule \vspace{-5pt}]

\pdfgentounicode=1

\newcommand{\resumeItem}[1]{
  \item\small{{#1 \vspace{-2pt}}}
}

\newcommand{\resumeSubheading}[4]{
  \vspace{-2pt}\item
    \begin{tabular*}{1.0\textwidth}[t]{l@{\extracolsep{\fill}}r}
      \textbf{#1} & \textbf{\small #2} \\
      \textit{\small#3} & \textit{\small #4} \\
    \end{tabular*}\vspace{-7pt}
}

\newcommand{\resumeProjectHeading}[2]{
    \item
    \begin{tabular*}{0.97\textwidth}{l@{\extracolsep{\fill}}r}
      \small#1 & \textbf{\small #2}\\
    \end{tabular*}\vspace{-7pt}
}

\renewcommand\labelitemi{$\vcenter{\hbox{\tiny$\bullet$}}$}
\renewcommand\labelitemii{$\vcenter{\hbox{\tiny$\bullet$}}$}

\newcommand{\resumeSubHeadingListStart}{\begin{itemize}[leftmargin=0.0in, label={}]}
\newcommand{\resumeSubHeadingListEnd}{\end{itemize}}
\newcommand{\resumeItemListStart}{\begin{itemize}}
\newcommand{\resumeItemListEnd}{\end{itemize}\vspace{-5pt}}

% ---------------------------------------------------------------------
% \cvheader{positioning line}{focus line}
% ---------------------------------------------------------------------
\newcommand{\cvheader}[2]{%
\begin{center}
    {\Huge \scshape Franck Deba} \\[0.4em]
    {\small \textit{#1}} \\[0.25em]
    {\small #2}
\end{center}
\vspace{0.5pt}
\small
\raisebox{-0.1\height}{\faPhone}~+33 (0)6 05 77 47 51 \hspace{1em}
\raisebox{-0.2\height}{\faEnvelope}~\href{mailto:debafranck@gmail.com}{\underline{debafranck@gmail.com}} \hspace{1em}
\raisebox{-0.2\height}{\faLinkedin}~\href{https://linkedin.com/in/deba-franck}{\underline{deba-franck}} \hspace{1em}
\raisebox{-0.2\height}{\faGithub}~\href{https://github.com/usenameFD}{\underline{usenameFD}} \hspace{1em}
\raisebox{-0.2\height}{\faGlobe}~\href{https://usenamefd.github.io/}{\underline{usenamefd.github.io}}
\vspace{-8pt}
}

%  Do not put content here — only packages, formatting and macros.
% =====================================================================
\documentclass[letterpaper,11pt]{article}

\usepackage{latexsym}
\usepackage[empty]{fullpage}
\usepackage{titlesec}
\usepackage{marvosym}
\usepackage[usenames,dvipsnames]{color}
\usepackage{verbatim}
\usepackage{enumitem}
\usepackage[hidelinks]{hyperref}
\usepackage{fancyhdr}
\usepackage[english]{babel}
\usepackage{tabularx}
\usepackage{fontawesome5}
\usepackage{multicol}
\usepackage{graphicx}
\setlength{\multicolsep}{-3.0pt}
\setlength{\columnsep}{-1pt}
\input{glyphtounicode}
\usepackage{paracol}

\pagestyle{fancy}
\fancyhf{}
\fancyfoot{}
\renewcommand{\headrulewidth}{0pt}
\renewcommand{\footrulewidth}{0pt}

\addtolength{\oddsidemargin}{-0.5in}
\addtolength{\evensidemargin}{-0.5in}
\addtolength{\textwidth}{0.9in}
\addtolength{\topmargin}{-.3in}
\addtolength{\textheight}{0.9in}
\hypersetup{
    colorlinks=true,
    linkcolor=blue,
    urlcolor=blue,
    filecolor=blue,
    pdfpagemode=FullScreen
}
\urlstyle{same}

\raggedbottom
\raggedright
\setlength{\tabcolsep}{0in}

\titleformat{\section}{
  \vspace{-4pt}\scshape\raggedright\large\bfseries
}{}{0em}{}[\color{black}\titlerule \vspace{-5pt}]

\pdfgentounicode=1

\newcommand{\resumeItem}[1]{
  \item\small{{#1 \vspace{-2pt}}}
}

\newcommand{\resumeSubheading}[4]{
  \vspace{-2pt}\item
    \begin{tabular*}{1.0\textwidth}[t]{l@{\extracolsep{\fill}}r}
      \textbf{#1} & \textbf{\small #2} \\
      \textit{\small#3} & \textit{\small #4} \\
    \end{tabular*}\vspace{-7pt}
}

\newcommand{\resumeProjectHeading}[2]{
    \item
    \begin{tabular*}{0.97\textwidth}{l@{\extracolsep{\fill}}r}
      \small#1 & \textbf{\small #2}\\
    \end{tabular*}\vspace{-7pt}
}

\renewcommand\labelitemi{$\vcenter{\hbox{\tiny$\bullet$}}$}
\renewcommand\labelitemii{$\vcenter{\hbox{\tiny$\bullet$}}$}

\newcommand{\resumeSubHeadingListStart}{\begin{itemize}[leftmargin=0.0in, label={}]}
\newcommand{\resumeSubHeadingListEnd}{\end{itemize}}
\newcommand{\resumeItemListStart}{\begin{itemize}}
\newcommand{\resumeItemListEnd}{\end{itemize}\vspace{-5pt}}

% ---------------------------------------------------------------------
% \cvheader{positioning line}{focus line}
% ---------------------------------------------------------------------
\newcommand{\cvheader}[2]{%
\begin{center}
    {\Huge \scshape Franck Deba} \\[0.4em]
    {\small \textit{#1}} \\[0.25em]
    {\small #2}
\end{center}
\vspace{0.5pt}
\small
\raisebox{-0.1\height}{\faPhone}~+33 (0)6 05 77 47 51 \hspace{1em}
\raisebox{-0.2\height}{\faEnvelope}~\href{mailto:debafranck@gmail.com}{\underline{debafranck@gmail.com}} \hspace{1em}
\raisebox{-0.2\height}{\faLinkedin}~\href{https://linkedin.com/in/deba-franck}{\underline{deba-franck}} \hspace{1em}
\raisebox{-0.2\height}{\faGithub}~\href{https://github.com/usenameFD}{\underline{usenameFD}} \hspace{1em}
\raisebox{-0.2\height}{\faGlobe}~\href{https://usenamefd.github.io/}{\underline{usenamefd.github.io}}
\vspace{-8pt}
}


\begin{document}

\cvheader{Quantitative Analyst --- Asset Management}{Portfolio construction $\cdot$ Index replication $\cdot$ Liquidity risk $\cdot$ ESG integration $\cdot$ Asset pricing}

%-----------EDUCATION-----------
\section{Education}
\resumeSubHeadingListStart
  \resumeSubheading
    {Master M2MO --- Quantitative Finance and Data Science}{Sep 2025 -- Present}
    {Universit\'e Paris Cit\'e}{Paris, France}
  \vspace{3pt}
  \resumeSubheading
    {MSc in Data Science and Risk Management}{Aug 2023 -- Dec 2025}
    {ENSAI}{Rennes, France}
  \vspace{3pt}
  \resumeSubheading
    {MSc in Statistics applied to Economics}{Sep 2021 -- Aug 2023}
    {ISSEA -- CEMAC}{Yaound\'e, Cameroon}
  \vspace{3pt}
  \resumeSubheading
    {Bachelor's Degree in Mathematics}{Sep 2017 -- Aug 2020}
    {Higher Teacher Training College \& University of Yaound\'e I}{Yaound\'e, Cameroon}
\resumeSubHeadingListEnd

%-----------PORTFOLIO PROJECTS-----------
\section{Portfolio \& Asset Management Projects}
\resumeSubHeadingListStart

  \resumeProjectHeading
    {\href{https://usenamefd.github.io/projects/risques-asset-management.html}{\textbf{Integrated Risk Framework for an Asset Manager}} $|$ \emph{Python}} {2025}
    \resumeItemListStart
      \resumeItem{Mean--variance portfolio construction under a target-return constraint, with annualised return/volatility estimation and correlation analysis}
      \resumeItem{Liquidity run-off profiles under four configurations (normal vs.\ stressed, with and without behavioural deformation), used to derive an \emph{optimal portfolio size} consistent with redemption constraints}
      \resumeItem{Fixed-income sleeve: bond valuation with recovery, rate sensitivity, Z-VaR vs.\ full-repricing VaR, and joint rate/credit-spread VaR via the credit triangle}
      \resumeItem{Stress testing of portfolio value and liquidity characteristics}
    \resumeItemListEnd

  \resumeProjectHeading
    {\href{https://usenamefd.github.io/projects/courbes-de-taux.html}{\textbf{Index Replication \& Interest Rate Modelling (CAC 40)}} $|$ \emph{Python, Dash}} {2025}
    \resumeItemListStart
      \resumeItem{Built a long-only replicating portfolio by tracking-error minimisation with position caps; out-of-sample backtest on tracking error, information ratio, beta and correlation}
      \resumeItem{Zero-coupon curve construction by bootstrapping and spline interpolation; Hull--White calibration and Monte Carlo pricing of rate derivatives}
      \resumeItem{Delivered through a Dash interface for allocation and replication analysis}
    \resumeItemListEnd

  \resumeProjectHeading
    {\href{https://usenamefd.github.io/projects/univers-investissement-esg.html}{\textbf{ESG Investment Universe Construction}} $|$ \emph{Python, pandas}} {2025}
    \resumeItemListStart
      \resumeItem{Two-stage methodology: normative exclusion (UN Global Compact) and sectoral exclusion of controversial activities, then best-in-class selection}
      \resumeItem{Median E, S and G scores of a reference ESG index used as benchmarks; simultaneous filtering on all three pillars, preventing compensation between dimensions}
      \resumeItem{Tracked coverage ratio at each filter to quantify the diversification cost of the ESG constraint}
    \resumeItemListEnd

  \resumeProjectHeading
    {\href{https://usenamefd.github.io/projects/bibliotheque-pricing-dash.html}{\textbf{Asset Pricing \& Allocation Toolkit}} $|$ \emph{Python, Dash}} {2025}
    \resumeItemListStart
      \resumeItem{Markowitz optimisation and tangency portfolio from empirical covariance, with correlation heatmaps}
      \resumeItem{Option pricing with Greeks, Asian and barrier payoffs; bond, FRA and swap valuation from a reconstructed risk-free curve}
    \resumeItemListEnd

  \resumeProjectHeading
    {\href{https://usenamefd.github.io/projects/creditvar-copules.html}{\textbf{Multi-Asset Dependence \& CreditVaR}} $|$ \emph{Python, Monte Carlo}} {2025}
    \resumeItemListStart
      \resumeItem{Monte Carlo CreditVaR at 99\% for a portfolio of bank receivables, from market-implied default probabilities and estimated recovery distributions}
      \resumeItem{Dependence modelled with elliptical (Gaussian, Student) and Archimedean (Gumbel, Clayton, Frank) copulas; dependogram evidence of upper-tail dependence that a Gaussian copula cannot reproduce}
    \resumeItemListEnd

  \resumeProjectHeading
    {\href{https://usenamefd.github.io/projects/econometrie-financiere-vecm.html}{\textbf{Time Series \& Relative-Value Modelling}} $|$ \emph{Python, R, Dash}} {2024}
    \resumeItemListStart
      \resumeItem{Cointegration testing and VECM calibration exploited as a mean-reversion (pair trading) strategy}
      \resumeItem{VAR with FARIMA--GARCH volatility, LSTM and Prophet for multi-horizon equity forecasting, deployed via Dash}
    \resumeItemListEnd

\resumeSubHeadingListEnd

%-----------EXPERIENCE-----------
\section{Professional Experience}
\resumeSubHeadingListStart
  \resumeSubheading
    {\textbf{Quantitative Analyst (Intern)} $|$ \textbf{PwC}}{Apr -- Oct 2026}
    {Risk Line of Service}{Paris, France}
    \resumeItemListStart
      \resumeItem{Modelling of the Default Risk Charge (DRC) under the FRTB capital framework, on an equity portfolio of CAC 40 issuers}
    \resumeItemListEnd

  \resumeSubheading
    {\textbf{Quantitative Risk Analyst (Intern)} $|$ \textbf{Cr\'edit Agricole}}{Apr -- Sep 2025}
    {Risk Modeling Department}{Paris, France}
    \resumeItemListStart
      \resumeItem{Developed climate stress-testing methodologies using time series and Monte Carlo methods}
      \resumeItem{Built an automated backtesting tool in Python (Dash) for model validation}
      \resumeItem{Migrated models from SAS to a Python-compatible platform}
    \resumeItemListEnd

  \resumeSubheading
    {\textbf{Data Science Intern} $|$ \textbf{Groupe VYV -- Harmonie Mutuelle}}{Jun -- Aug 2024}
    {Assurance Pr\'evoyance}{Paris, France}
    \resumeItemListStart
      \resumeItem{Built insured-population profiles by k-means clustering and estimated expenditure by risk profile}
      \resumeItem{Quantified risk of prolonged treatment via survival analysis (Cox model)}
    \resumeItemListEnd
\resumeSubHeadingListEnd

%-----------SKILLS-----------
\section{Technical Skills}
\begin{itemize}[leftmargin=0.15in, label={}]
  \small{\item{
    \textbf{Languages:} Python (NumPy, SciPy, pandas, scikit-learn, PyTorch), R, SQL, SAS, C++ (basics) \\
    \textbf{Portfolio \& Risk:} Mean--variance optimisation, tangency portfolio, tracking-error minimisation, liquidity run-off, VaR / Expected Shortfall, copulas, extreme value theory \\
    \textbf{Frameworks:} FRTB / CRR3, Basel III, IRB, SFDR, UN Global Compact, SBTi \\
    \textbf{Tools:} Dash, Git, Linux, \LaTeX
  }}
\end{itemize}

%-----------CERTIFICATIONS-----------
\section{Certifications, Awards \& Languages}
\begin{itemize}[leftmargin=0.15in, label={}]
  \small{\item{
    \textbf{Certifications:} AMF (Autorit\'e des March\'es Financiers), Financial Markets (Yale), AI in Finance (World Bank) \\
    \textbf{Awards:} National Essay Competition Prize (2019), PK Fokam Prize for Science Students (2017) \\
    \textbf{Languages:} French (native), English (advanced, B2 --- TOEIC/TOEFL)
  }}
\end{itemize}

\end{document}
